\section{Problem Analysis}

Modern children face severe attention span challenges that make traditional programming education ineffective.
Standard programming courses require prolonged focus, abstract thinking, and patience for delayed results. 
Interface complexity becomes an obstacle instead of supporting the learning process. These requirements clash with how children consume information today. 

Speaking of existing tools, there are technical barriers. Children face Java requirements far above beginner level when working with Minecraft mods. 
Tools like ComputerCraft use Lua syntax that confuses young learners. 
Setup complexity takes time before anyone writes a single line of code, which stops learning early. 
Command blocks, an in-game programming interface, overwhelm beginners from the start. 
Parents or teachers must handle installation and configuration because children lack the technical knowledge.

As for pedagogical gaps, children learn abstract concepts without seeing corresponding actions in the game. 
Block-based tools hide text syntax and create barriers when transitioning to real code. 
Motivation drops when results feel distant or unclear. 
Available tools do not match how children interact with digital products today. 
Educational resources assume adult knowledge or prior programming experience.

Pain Points in Existing Solutions for Teaching Programming:

\begin{itemize}
    \item[\textbf{--}] \textbf{Scratch and Other Visual Languages:}
    Scratch lowers the entry barrier by removing syntax errors, but it operates in an abstract environment disconnected from the games children already play. 
    The visual block system hides real text-based programming syntax, which makes the transition to textual languages like Python or Java difficult \cite{noauthor_connecting_2024}. 
    As projects grow in complexity, block-based coding limits the types of programs learners can build \cite{guide_scratch_2023}. 
    Learners often struggle to connect their work to familiar gaming environments, causing engagement to drop \cite{adler_improving_nodate}.

    \item[\textbf{--}] \textbf{Minecraft Education Edition:}
    Minecraft Education Edition integrates learning into a familiar game, but access requires institutional licensing, limiting home use \cite{noauthor_licensing_nodate}. 
    The platform targets classroom instruction and emphasizes guided lessons over open-ended experimentation \cite{noauthor_pros_2025}. 
    Activities rely heavily on pre-built curricula, which reduces creative freedom and limits exploration. 
    Many children prefer the regular version of Minecraft where friends play, making the educational edition feel isolated from their social gaming environment \cite{sripan_gamified_2025}.

    \item[\textbf{--}] \textbf{ComputerCraft:}
    ComputerCraft introduces real programming concepts through Lua, but it requires text-based coding from the start. 
    Installation and setup demand technical knowledge often unavailable to parents. 
    Documentation assumes prior programming experience, which can overwhelm beginners. 
    Young learners face high complexity before achieving results, leading to frustration and early disengagement \cite{kazemitabaar_scaffolding_2023}.

    \item[\textbf{--}] \textbf{Other Common Approaches and Their Limitations:}
    Minecraft command blocks allow in-game automation, but dense syntax produces frequent errors that confuse beginners. 
    Redstone teaches logic indirectly, and debugging circuits is unclear, making reasoning about failures difficult. 
    Traditional programming languages such as Python or Java require setup and configuration before learners can see results; syntax errors in early stages often block progress and reduce motivation.
\end{itemize}


Thus, the problem statement is the following: \textbf{Children disengage from programming education when learning tools fail to provide immediate visual feedback and do not connect abstract coding concepts to activities that sustain their interest.}